% GENERATED FILE. Edit canonical lessons, then run npm run generate:books.
% canonical-source-hash: fnv1a64:130eec18c25b679d
% canonical-lessons: GE-C294-katastrophe, GE-C294-erdbeben, GE-C294-ueberschwemmung, GE-C294-hochwasser, GE-C294-einstuerzen

\chapter{Naturkatastrophen}
\label{ch:ge-naturkatastrophen}
\hlchaptermodality{294}

\begin{chapteropening}
\textbf{By the end of this chapter:} \emph{I can report a natural disaster: a disaster, an earthquake, a flood, high water and to collapse.}

Complete the last lesson of chapter 294: I can report a natural disaster: a disaster, an earthquake, a flood, high water and to collapse.
\end{chapteropening}

\section[die Katastrophe]{die Katastrophe — the disaster, the catastrophe}

\label{lesson:GE-C294-katastrophe}



\begin{quote}
Before the new one: say the German for to put out, then the German for the forest fire.
\end{quote}

\subsection*{You'll want to know: die Katastrophe}
\textbf{die Katastrophe} — \textquotedblleft{}the disaster, the catastrophe\textquotedblright{}.

\begin{quote}
Das Wetter war eine Katastrophe.
\end{quote}

\begin{quote}
\emph{The weather was a disaster.}
\end{quote}

Plural: \emph{die Katastrophen}. Colloquially, anything that went badly.

\begin{cousinweb}
\emph{Katastrophe} comes from Latin \emph{catastropha}, from Greek \emph{katastrophē} \textquotedblleft{}an overturning\textquotedblright{} (EtymDB), built on \emph{kata} \textquotedblleft{}down\textquotedblright{} and \emph{strephein} \textquotedblleft{}to turn\textquotedblright{}.
\end{cousinweb}

\begin{grammarlens}[title={What you've built}]
One more word for telling, explaining and describing at B1.
\end{grammarlens}

\subsection*{Your turn}
\begin{itemize}
\raggedright
  \item \emph{Say it:} \emph{die Katastrophe}
  \item \emph{Say it:} \emph{die Katastrophe}, once more
  \item \emph{Say it:} a sentence of your own with \emph{die Katastrophe}
  \item \emph{Recall:} say the German for to put out, then the German for the forest fire, then say \emph{die Katastrophe} again
\end{itemize}

\begin{tcolorbox}[breakable,colback=teal!4,colframe=teal!35!black,title={Before you move on}]
What does die Katastrophe mean? (\textquotedblleft{}The disaster, the catastrophe\textquotedblright{}.) Say it once more.
\end{tcolorbox}
\section[das Erdbeben]{das Erdbeben — the earthquake}

\label{lesson:GE-C294-erdbeben}



\begin{quote}
Before the new one: say the German for the forest fire, then the German for the disaster.
\end{quote}

\subsection*{You'll want to know: das Erdbeben}
\textbf{das Erdbeben} — \textquotedblleft{}the earthquake\textquotedblright{}.

\begin{quote}
Das Erdbeben dauerte nur zehn Sekunden.
\end{quote}

\begin{quote}
\emph{The earthquake lasted only ten seconds.}
\end{quote}

Plural: \emph{die Erdbeben}. \emph{Erde} 'earth' plus \emph{beben} 'to shake'.

\begin{cousinweb}
\emph{Erdbeben} joins \emph{Erde} \textquotedblleft{}earth\textquotedblright{} and \emph{beben} \textquotedblleft{}to tremble\textquotedblright{}, from Old High German \emph{bibēn}, Proto-Germanic \emph{bibjaną} \textquotedblleft{}to quake\textquotedblright{} (EtymDB).
\end{cousinweb}

\begin{grammarlens}[title={What you've built}]
One more word for telling, explaining and describing at B1.
\end{grammarlens}

\subsection*{Your turn}
\begin{itemize}
\raggedright
  \item \emph{Say it:} \emph{das Erdbeben}
  \item \emph{Say it:} \emph{das Erdbeben}, once more
  \item \emph{Say it:} a sentence of your own with \emph{das Erdbeben}
  \item \emph{Recall:} say the German for the forest fire, then the German for the disaster, then say \emph{das Erdbeben} again
\end{itemize}

\begin{tcolorbox}[breakable,colback=teal!4,colframe=teal!35!black,title={Before you move on}]
What does das Erdbeben mean? (\textquotedblleft{}The earthquake\textquotedblright{}.) Say it once more.
\end{tcolorbox}
\section[die Überschwemmung]{die Überschwemmung — the flood}

\label{lesson:GE-C294-ueberschwemmung}



\begin{quote}
Before the new one: say the German for the disaster, then the German for the earthquake.
\end{quote}

\subsection*{You'll want to know: die Überschwemmung}
\textbf{die Überschwemmung} — \textquotedblleft{}the flood\textquotedblright{}.

\begin{quote}
Nach dem Regen gab es eine Überschwemmung im Tal.
\end{quote}

\begin{quote}
\emph{After the rain there was a flood in the valley.}
\end{quote}

Plural: \emph{die Überschwemmungen}.

\begin{cousinweb}
\emph{Überschwemmung} is built on \emph{überschwemmen} \textquotedblleft{}to flood\textquotedblright{}, from \emph{schwemmen} \textquotedblleft{}to wash along\textquotedblright{}, the causative of \emph{schwimmen} \textquotedblleft{}to swim\textquotedblright{}; EtymDB lists no parent.
\end{cousinweb}

\begin{grammarlens}[title={What you've built}]
One more word for telling, explaining and describing at B1.
\end{grammarlens}

\subsection*{Your turn}
\begin{itemize}
\raggedright
  \item \emph{Say it:} \emph{die Überschwemmung}
  \item \emph{Say it:} \emph{die Überschwemmung}, once more
  \item \emph{Say it:} a sentence of your own with \emph{die Überschwemmung}
  \item \emph{Recall:} say the German for the disaster, then the German for the earthquake, then say \emph{die Überschwemmung} again
\end{itemize}

\begin{tcolorbox}[breakable,colback=teal!4,colframe=teal!35!black,title={Before you move on}]
What does die Überschwemmung mean? (\textquotedblleft{}The flood\textquotedblright{}.) Say it once more.
\end{tcolorbox}
\section[das Hochwasser]{das Hochwasser — the high water, a river in flood}

\label{lesson:GE-C294-hochwasser}



\begin{quote}
Before the new one: say the German for the earthquake, then the German for the flood.
\end{quote}

\subsection*{You'll want to know: das Hochwasser}
\textbf{das Hochwasser} — \textquotedblleft{}the high water, a river in flood\textquotedblright{}.

\begin{quote}
Das Hochwasser kam bis an unsere Tür.
\end{quote}

\begin{quote}
\emph{The high water came up to our door.}
\end{quote}

Plural: \emph{die Hochwasser}. For rivers that rise.

\begin{cousinweb}
\emph{Hochwasser} joins \emph{hoch} \textquotedblleft{}high\textquotedblright{}, Old High German \emph{hōh} (EtymDB), and \emph{Wasser} \textquotedblleft{}water\textquotedblright{}, both inherited Germanic words.
\end{cousinweb}

\begin{grammarlens}[title={What you've built}]
One more word for telling, explaining and describing at B1.
\end{grammarlens}

\subsection*{Your turn}
\begin{itemize}
\raggedright
  \item \emph{Say it:} \emph{das Hochwasser}
  \item \emph{Say it:} \emph{das Hochwasser}, once more
  \item \emph{Say it:} a sentence of your own with \emph{das Hochwasser}
  \item \emph{Recall:} say the German for the earthquake, then the German for the flood, then say \emph{das Hochwasser} again
\end{itemize}

\begin{tcolorbox}[breakable,colback=teal!4,colframe=teal!35!black,title={Before you move on}]
What does das Hochwasser mean? (\textquotedblleft{}The high water, a river in flood\textquotedblright{}.) Say it once more.
\end{tcolorbox}
\section[einstürzen]{einstürzen — to collapse, to fall in}

\label{lesson:GE-C294-einstuerzen}



\begin{quote}
Before the new one: say the German for the flood, then the German for the high water.
\end{quote}

\subsection*{You'll want to know: einstürzen}
\textbf{einstürzen} — \textquotedblleft{}to collapse, to fall in\textquotedblright{}.

\begin{quote}
Die alte Brücke stürzte ein.
\end{quote}

\begin{quote}
\emph{The old bridge collapsed.}
\end{quote}

Separable, Perfekt with \emph{sein}: \emph{Das Haus ist eingestürzt}.

\begin{cousinweb}
\emph{einstürzen} joins \emph{ein} \textquotedblleft{}in\textquotedblright{} and \emph{stürzen} \textquotedblleft{}to plunge\textquotedblright{}, from Old High German \emph{sturzen} (EtymDB).
\end{cousinweb}

\begin{grammarlens}[title={What you've built}]
One more word for telling, explaining and describing at B1.
\end{grammarlens}

\subsection*{Your turn}
\begin{itemize}
\raggedright
  \item \emph{Say it:} \emph{einstürzen}
  \item \emph{Say it:} \emph{einstürzen}, once more
  \item \emph{Say it:} a sentence of your own with \emph{einstürzen}
  \item \emph{Recall:} say the German for the flood, then the German for the high water, then say \emph{einstürzen} again
\end{itemize}

\begin{tcolorbox}[breakable,colback=teal!4,colframe=teal!35!black,title={Before you move on}]
What does einstürzen mean? (\textquotedblleft{}To collapse, to fall in\textquotedblright{}.) Say it once more.
\end{tcolorbox}
